\subsection{Projective Cover of a Module}

DEFINITION 3. --- Let A be a ring, and let M be an A-module. A projective cover of M is a pair \((P, u)\) , where P is a projective A-module and u a surjective homomorphism from P to M such that we have \(u(P') \neq M\) for every proper A-submodule \(P'\) of P.

Remarks. --- 1) For every projective A-module M, the pair \((M,1_{M})\) is a projective cover of M.

\begin{enumerate}
\def\labelenumi{\arabic{enumi})}
\setcounter{enumi}{1}
\tightlist
\item
  Suppose that \((\mathrm{P}, u)\) is a projective cover of the A-module M. Let \((x_{i})_{i \in \mathrm{I}}\) be a family of elements of P, and let \(P'\) be the submodule of P that it generates; then \(u(\mathrm{P}')\) is generated by the family \((u(x_{i}))_{i \in \mathrm{I}}\) . Consequently, the family \((x_{i})_{i \in \mathrm{I}}\) generates the A-module P if and only if the family \((u(x_{i}))_{i \in \mathrm{I}}\) generates the A-module M. In particular, P is finitely generated if and only if M is finitely generated.
\end{enumerate}

PROPOSITION 8. --- Let M and M' be A-modules, (P, u) and (P', u') projective covers of M and M', respectively, and g : M → M' an A-linear mapping.

\begin{enumerate}
\def\labelenumi{\alph{enumi})}
\item
  There exists an A-linear mapping \(f: \mathrm{P} \to \mathrm{P}'\) such that \(u' \circ f = g \circ u\) .
\item
  Let f be such a mapping. If g is surjective (resp. bijective), then f is surjective (resp. bijective). If g is injective and its image is a direct factor of \(M'\) , then f is injective and its image is a direct factor of \(P'\) .
\end{enumerate}

By assumption, the A-module P is projective and the mapping \(u'\) is surjective. Consequently, there exists an A-linear mapping \(f: P \to P'\) such that \(g \circ u = u' \circ f\) (II, §2, No.~2, p.~231, Proposition 4). Assertion a) follows.

Let \(f\) be a mapping as in a). Suppose that \(g\) is surjective. Since \(u\) is surjective, we have \(M' = g(u(P)) = u'(f(P))\) . Since \((P', u')\) is a projective cover of \(M'\) , we have \(f(P) = P'\) , so \(f\) is surjective. By loc. cit., the kernel of \(f\) admits a supplementary submodule \(P_1\) , so \(f(P_1) = P'\) . Now suppose that \(g\) is bijective. We have \(g(u(P_1)) = u'(f(P_1)) = u'(P') = M'\) , hence \(u(P_1) = M\) . Since \((P, u)\) is a projective cover of \(M\) , we have \(P_1 = P\) , hence \(\text{Ker}(f) = 0\) . So \(f\) is injective; since we already know that \(f\) is surjective, \(f\) is bijective.

To conclude, suppose that \(g\) is injective and that its image is a direct factor of \(\mathrm{M}'\) . Then there exists an A-linear mapping \(g': \mathrm{M}' \to \mathrm{M}\) such that \(g' \circ g = 1_{\mathrm{M}}\) . By a), there exists an A-linear mapping \(f': \mathrm{P}' \to \mathrm{P}\) such that \(u \circ f' = g' \circ u'\) . We have \(u \circ (f' \circ f) = g' \circ u' \circ f = (g' \circ g) \circ u\) ; by the previous paragraph, the mapping \(f' \circ f\) is bijective. Denote the converse bijection by h; we have \((h \circ f') \circ f = 1_{\mathrm{P}}\) , so f is injective, and its image is a direct factor of \(P'\) (II, §1, No.~9, p.~212, Corollary 2).

COROLLARY 1. --- Let M be an A-module. Let \((P, u)\) and \((P', u')\) be projective covers of M. There exists an isomorphism f from P to \(P'\) such that \(u = u' \circ f\) .

\pandocbounded{\includegraphics[keepaspectratio]{images/part001_525bb1e1.jpg}}

Note that f is not necessarily uniquely determined by the relation \(u = u' \circ f\) (VIII, p.~170, Exercise 21).

COROLLARY 2. --- Let \((\mathrm{P}, u)\) be a projective cover of the A-module M. If Q is a projective A-module and \(g : Q \to M\) a surjective linear mapping, then there exists a surjective linear mapping \(f : Q \to P\) such that \(g = u \circ f\) .

PROPOSITION 9. --- Let M be an A-module and \((P, u)\) a projective cover of M. Denote the radical of the ring A by r. The homomorphism \(\overline{u}: P/rP \to M/rM\) deduced from u by passing to the quotients is an isomorphism.

The homomorphism \(u\) is surjective by definition, so \(\overline{u}\) is surjective. Denote the kernel of \(u\) by \(N\) . We have \(u^{-1}(\mathfrak{rM}) = N + \mathfrak{rP}\) . Let us prove that we have \(N \subset \mathfrak{rP}\) , which implies the injectivity of \(\overline{u}\) . For every maximal submodule \(P'\) of \(P\) , we have \(u(P') \neq M\) , hence \(P' + N \neq P\) ; since \(P'\) is maximal, we have \(N \subset P'\) . The submodule \(N\) of \(P\) is therefore contained in the radical of \(P\) . Now, the latter is equal to \(\mathfrak{rP}\) by Proposition 6 of VIII, p.~158.

COROLLARY. --- If an A-module M has a projective cover, then the A/r-module M/rM is projective.

Indeed, if \((P,u)\) is a projective cover of M, then the \((A/\mathfrak{r})\) -module M/rM is isomorphic to P/rP (Proposition 9). Since the A-module P is projective, the \((A/\mathfrak{r})\) -module P/rP is also projective.

Remarks. --- 3) Suppose that the ring A is without radical. By Proposition 9, \((P, u)\) is a projective cover of an A-module M if and only if u is an isomorphism. Hence, projective covers can only exist for projective A-modules.

The ring Z is without radical (VIII, p.~154, Example 3). Let \(n \geqslant 2\) be an integer. The Z-module Z/nZ is not projective and therefore does not admit a projective cover.

\begin{enumerate}
\def\labelenumi{\arabic{enumi})}
\setcounter{enumi}{3}
\tightlist
\item
  Suppose that every finitely generated A-module has a projective cover; then the quotient \(A'\) of A by its radical is semisimple. Indeed, every finitely generated module over the ring \(A'\) is projective by the corollary. In particular, for every left ideal \(\mathfrak{a}\) of \(A'\) , the \(A'\) -module \(A_s'/\mathfrak{a}\) is projective. Our assertion then follows from Proposition 4 of VIII, p.~138.
\end{enumerate}

We can give examples of a commutative ring A for which A/r is semisimple and of a finitely generated A-module M that does not admit a projective cover (VIII, p.~170, Exercise 22).

PROPOSITION 10. --- Let M be an A-module, P a projective A-module, and \(u : P \to M\) a linear mapping. Let a be a two-sided ideal of A. Suppose that the linear mapping \(\overline{u} : P / aP \to M / aM\) deduced from u by passing to the quotients is bijective and that one of the following conditions are satisfied:

\begin{enumerate}
\def\labelenumi{(\roman{enumi})}
\item
  The A-modules M and P are finitely generated, and a is contained in the radical of A.
\item
  The ideal \(\mathfrak{a}\) is nilpotent.
\end{enumerate}

Then (P, u) is a projective cover of M.

Under the assumptions, the homomorphism \(u\) is surjective (VIII, p.~159, Corollary 2) and its kernel \(N\) is contained in \(\mathfrak{a}P\) . Let \(P'\) be a proper submodule of \(P\) . By Nakayama's lemma (VIII, p.~158, Theorem 2), we have \(P' + \mathfrak{a}P \neq P\) and therefore \(u(P') \neq M\) . So \((P, u)\) is a projective cover of \(M\) .

COROLLARY 1. --- Let P be a projective A-module. Suppose that P is finitely generated or that the radical r of A is a nilpotent two-sided ideal. Denote the canonical mapping from P to P/rP by u. Then (P, u) is a projective cover of P/rP.

COROLLARY 2. --- Let a be a two-sided ideal of A and M an A-module such that the (A/a)-module M/AM is free. Suppose that one of the following conditions is satisfied:

\begin{enumerate}
\def\labelenumi{(\roman{enumi})}
\item
  The module M is finitely generated, and \(\mathfrak{a}\) is contained in the radical of A.
\item
  The ideal \(\mathfrak{a}\) is nilpotent.
\end{enumerate}

Then M has a projective cover.

More precisely, let P be a free A-module, \((e_{i})_{i\in\mathrm{I}}\) a basis of P, and let \(u:P\to M\) be a homomorphism such that the canonical images of the elements \(u(e_{i})\) of M/aM form a basis of the (A/a)-module M/aM. Then (P,u) is a projective cover of M.

The \((\mathrm{A}/\mathfrak{a})\) -module P/aP is free, and the homomorphism \(\overline{u}\) from P/aP to M/aM deduced from u by passing to the quotients transforms a basis of P/aP into a basis of M/aM, hence is bijective.

If a is nilpotent, then it suffices to apply Proposition 10. Now suppose that the ring A is nonzero and that the module M is finitely generated; then so is the \((A/a)\) -module M/aM, and consequently also the \((A/a)\) -module P/aP. Every basis of P/aP is then finite. It follows that the set I is finite and that the A-module P is finitely generated. We then apply Proposition 10 again.

COROLLARY 3. --- Every finitely generated module over a local ring has a projective cover.

Let A be a local ring and r its radical. It is a two-sided ideal of A (VIII, p.~155, Proposition 5, a)), and the ring A/r is a field (VIII, p.~157, Corollary 4). If M is an A-module, then M/rM is a vector space over the field A/r, hence is a free (A/r)-module. It then suffices to apply Corollary 2.

Remark 5. --- Let A be a local ring and r its radical. Let M be a finitely generated A-module, P a finitely generated projective A-module, and \(u : P \to M\) a homomorphism. By Corollary 6 of VIII, p.~36, the A-module P is free. Choose a basis \((e_{i})_{i \in I}\) of P. Set \(x_{i} = u(e_{i})\) , and denote the canonical image of \(x_{i}\) in M/rM by \(\overline{x}_{i}\) . The following properties are equivalent:

\begin{enumerate}
\def\labelenumi{(\roman{enumi})}
\item
  The pair \((\mathrm{P},u)\) is a projective cover of M.
\item
  The family \((x_{i})_{i\in \mathrm{I}}\) is a minimal generating family of the A-module M.
\item
  The family \((\overline{x}_{i})_{i\in\mathrm{I}}\) is a basis of the vector space M/rM over the field A/r.
\end{enumerate}

We have (i) \(\Rightarrow\) (ii) by Remark 2 of VIII, p.~161 and (iii) \(\Rightarrow\) (i) by Corollary 2. Furthermore, if the family \((x_{i})\) is a minimal generating family of the A-module M, then the family \((\overline{x}_{i})\) is a minimal generating family, that is, a basis, of the vector space M/rM over A/r (VIII, p.~158, Corollary 1).

PROPOSITION 11. --- Let A be a ring and a nilpotent two-sided ideal of A. Let M be a projective A/a-module. There exist a projective A-module P and a surjective A-linear mapping \(u: \mathrm{P} \to \mathrm{M}\) with kernel aP.

Such a pair \((P,u)\) is a projective cover of M viewed as an A-module.

There exists an A/ \(\alpha\) -module \(M'\) such that \(M \oplus M'\) is a free A/ \(\alpha\) -module. Choose a free A-module L and a surjective A-linear mapping \(v: L \to M \oplus M'\) with kernel aL. By Corollary 1 of Proposition 7 (VIII, p.~159), there exists a direct sum decomposition \(L = P \oplus P'\) such that \(v(P) = M\) and \(v(P') = M'\) . The A-module P is projective, and the A-linear mapping u from P to M that coincides with v on P is surjective, with kernel aP. The first assertion follows.

Let P be a projective A-module and u a homomorphism from P to M with kernel aP. By Proposition 10, the pair \((P,u)\) is a projective cover of M.

COROLLARY. --- Let P and P' be projective A-modules. If the modules P/ \(\alpha\) P and P'/ \(\alpha\) P' are isomorphic, then P and P' are isomorphic.

Since P/aP and \(P^{\prime}/aP^{\prime}\) are projective, the corollary follows from Proposition 11 and the uniqueness of the projective cover (VIII, p.~162, Corollary 1).

